%! Author = sriadityavedantam
%! Date = 4/8/26

\section{Regular Differential 1-Forms}

Let $X$ be a quasi-projective variety.
We define regular differential $1$-forms on $X$.
Let $f$ be a regular function on a neighbourhood $U$ of a point $x \in X$, so $f \in \c[U]$.
The differential of $f$ at $x$ is denoted
\[
    d_x f \in T_x^\upstar,
\]
the cotangent space at $x$.
If $X = \aa^n$ with coordinates $t_1,\ldots,t_n$, then $T_x \cong \c^n$.
For $f = t_i$, we have
\[
    d_x t_i : \c^n \to \c,
\]
which is the projection onto the $i$-th coordinate.
More generally, for $f \in \c[t_1,\ldots,t_n]$,
\[
    d_x f : \c^n \to \c,
    \quad
    (u_1,\ldots,u_n) \mapsto \sum_{i=1}^n \left.\frac{\partial f}{\partial t_i}\right|_x u_i.
\]

If $X$ is not affine, we work locally on affine neighbourhoods.
Define
\[
    \Phi(X)
    = \left\{\phi : X \to \coprod_{x \in X} T_x^\upstar \;\middle|\; \phi(x) \in T_x^\upstar \right\}.
\]

\begin{definition}[Regular Differential 1-Form]
    A regular $1$-form on $X$ is a function $\phi \in \Phi(X)$ such that for every $x \in X$, there exists a neighbourhood $U$ of $x$ with
    \[
        \phi|_U = \sum_i f_i\, d g_i,
    \]
    where $f_i, g_i \in \c[U]$.
\end{definition}

Equivalently, using sheaf language, a regular $1$-form is a global section of the cotangent sheaf.
Explicitly,
\[
    \sum_i f_i\, d g_i : U \to \coprod_{x \in U} T_x^\upstar,
    \quad
    x \mapsto \sum_i f_i(x)\, d_x g_i.
\]

\begin{remark}
    Let $\Omega[X]$ denote the set of all regular $1$-forms on $X$.
    Then $\Omega[X]$ is a $\c[X]$-module.
\end{remark}

We have the natural differential map
\[
    \begin{aligned}
        d : \c[X] &\to \Omega[X]\\
        f &\mapsto (x \mapsto d_x f),
    \end{aligned}
\]
which satisfies
\[
    d(f+g) = df + dg,
    \quad
    d(fg) = f\,dg + g\,df.
\]

\begin{theorem*}{
    If $X$ is projective, then $\Omega[X]$ is a finite-dimensional $\c$-vector space.
}
\end{theorem*}

Since $X$ is projective, $\c[X] = \c$, so $\Omega[X]$ is naturally a $\c$-vector space.
Define
\[
    \dim_\c \Omega[X] \coloneqq h^{1,0}.
\]

If $X$ is a smooth projective curve, then $h^{1,0}$ is the geometric genus of $X$.